\begin{examplebox}
  \textbf{\textcolor{CExample}{例 1（入门：从原函数的对称轴找导函数的对称中心）}}\par\nopagebreak[4]

  \textbf{题目：} $f(x)=(x-2)^2+1$。不必画出导函数的图像，判断它的对称中心，并核对 $f'(1)$、$f'(3)$。

  \textbf{解题步骤：}
  \begin{description}[style=nextline,leftmargin=2.8em,labelsep=0.5em,itemsep=0.3em]
    \item[\textbf{第一步（看原图）：}] $f$ 关于 $x=2$ 对称，左右等距点 $1$ 和 $3$ 的函数值都是 $2$。
    \item[\textbf{第二步（直接用结论）：}] $a=2$，所以 $f'$ 关于 $(2,0)$ 中心对称，$f'(1)=-f'(3)$。
    \item[\textbf{第三步（验算）：}] $f'(x)=2(x-2)$，于是 $f'(1)=-2$，$f'(3)=2$，完全符合。
  \end{description}
  \textbf{关键结论：}看见原函数的\textbf{对称轴 $x=a$}，立即联想到导函数的\textbf{对称中心 $(a,0)$}。
  \begin{center}
  \begin{tikzpicture}[font=\small,>=stealth,line join=round]
\draw[draw=teal!45!black,fill=green!3,rounded corners=3pt] (0,0) rectangle (13.6,6.5);
\node at (6.8,5.95) {同一对横坐标：原函数同高，导函数的值互为相反数};
\node at (3.5,5.15) {$f(x)=(x-2)^2+1$};
\node at (10.3,5.15) {$f^{\prime}(x)=2(x-2)$};
\begin{scope}[shift={(1.2,1.4)},x=0.95cm,y=0.55cm]
\draw[gray,->] (-0.2,0)--(4.4,0) node[right] {$x$};
\draw[gray,->] (0,-0.3)--(0,5.5);
\draw[teal!75!black,very thick,domain=0.3:3.7,samples=100] plot(\x,{(\x-2)^2+1});
\draw[dashed,teal!65!black] (2,0)--(2,5);
\node[above] at (2,5.1) {$x=2$};
\draw[dotted,gray] (1,2)--(3,2);
\fill[teal!75!black] (1,2) circle(2.5pt);
\fill[teal!75!black] (3,2) circle(2.5pt);
\node[above] at (1,2.8) {$2$};\node[above] at (3,2.8) {$2$};
\draw[gray] (1,0+.08)--(1,0-.08);\node[below] at (1,0-.15) {$1$};
\draw[gray] (2,0+.08)--(2,0-.08);\node[below] at (2,0-.15) {$2$};
\draw[gray] (3,0+.08)--(3,0-.08);\node[below] at (3,0-.15) {$3$};
\end{scope}
\begin{scope}[shift={(8.2,2.8)},x=0.95cm,y=0.46cm]
\draw[gray,->] (-0.2,0)--(4.4,0) node[right] {$x$};
\draw[gray,->] (0,-3.8)--(0,3.8);
\draw[blue!70!black,very thick,domain=0.3:3.7,samples=100] plot(\x,{2*(\x-2)});
\draw[dotted,gray] (0,-2)--(1,-2) (0,2)--(3,2);
\fill[blue!70!black] (1,-2) circle(2.5pt);
\fill[blue!70!black] (3,2) circle(2.5pt);
\fill[blue!70!black] (2,0) circle(2.5pt);
\node[left] at (0,-2) {$-2$};\node[left] at (0,2) {$2$};
\node[above left] at (1.9,.45) {$(2,0)$};
\draw[gray] (1,0+.08)--(1,0-.08);\node[below] at (1,0-.15) {$1$};
\draw[gray] (3,0+.08)--(3,0-.08);\node[below] at (3,0-.15) {$3$};
\end{scope}
\node at (6.8,.45) {$1,3$ 的中点是 $2$：左图看轴，右图看中心 $(2,0)$。};
\end{tikzpicture}


  \end{center}

  \medskip
  \textbf{\textcolor{CExample}{例 2（反推：只看导函数的对称轴）}}\par\nopagebreak[4]

  \textbf{题目：}已知
  \[
    f'(x)=3(x-1)^2+2,\qquad f(1)=7.
  \]
  不必先求 $f(x)$，判断原函数的对称中心，并求 $f(0)+f(2)$。

  \textbf{解题步骤：}
  \begin{description}[style=nextline,leftmargin=2.8em,labelsep=0.5em,itemsep=0.3em]
    \item[\textbf{第一步（观察导数）：}] $f'(x)$ 关于直线 $x=1$ 对称，所以 $f'(1+t)=f'(1-t)$。
    \item[\textbf{第二步（反推原函数）：}] 由第二组结论，$f$ 关于点 $(1,f(1))=(1,7)$ 中心对称。
    \item[\textbf{第三步（左右配对）：}] 横坐标 $0$ 与 $2$ 的中点是 $1$，所以 $f(0)+f(2)=2f(1)=\boxed{14}$。
  \end{description}
  \textbf{关键结论：}导函数轴对称时，原函数的\textbf{中心纵坐标一定是 $f(a)$}，不能随手写成 $0$。
  \begin{center}
  \begin{tikzpicture}[font=\small,>=stealth,line join=round]
\draw[draw=teal!45!black,fill=green!3,rounded corners=3pt] (0,0) rectangle (13.6,6.8);
\node at (6.8,6.25) {导函数关于 $x=1$ 对称，原函数的中心高度由 $f(1)=7$ 确定};
\node at (3.5,5.55) {$f^{\prime}(x)=3(x-1)^2+2$};
\node at (10.1,5.55) {$f(x)=(x-1)^3+2(x-1)+7$};
\begin{scope}[shift={(1.6,1.4)},x=1.15cm,y=0.4cm]
\draw[gray,->] (-0.2,0)--(3.6,0) node[right] {$x$};
\draw[gray,->] (0,-0.3)--(0,9.5);
\draw[blue!70!black,very thick,domain=-0.2:2.2,samples=100] plot(\x,{3*(\x-1)^2+2});
\draw[blue!65!black,dashed] (1,0)--(1,8.7);
\fill[blue!70!black] (0,5) circle(2.5pt);
\fill[blue!70!black] (2,5) circle(2.5pt);
\draw[dotted,gray] (0,5)--(2,5);
\node[above] at (0,5.7) {$5$};\node[above] at (2,5.7) {$5$};
\draw[gray] (0,0+.08)--(0,0-.08);\node[below] at (0,0-.15) {$0$};
\draw[gray] (1,0+.08)--(1,0-.08);\node[below] at (1,0-.15) {$1$};
\draw[gray] (2,0+.08)--(2,0-.08);\node[below] at (2,0-.15) {$2$};
\end{scope}
\begin{scope}[shift={(8.9,1.4)},x=1.15cm,y=0.32cm]
\draw[gray,->] (-0.35,0)--(3.4,0) node[right] {$x$};
\draw[gray,->] (0,-0.3)--(0,11.7);
\draw[teal!75!black,very thick,domain=-0.2:2.2,samples=100] plot(\x,{(\x-1)^3+2*(\x-1)+7});
\fill[teal!75!black] (0,4) circle(2.5pt);
\fill[blue!70!black] (1,7) circle(2.5pt);
\fill[teal!75!black] (2,10) circle(2.5pt);
\draw[dotted,gray] (0,4)--(2,10);
\node[left] at (-.12,4) {$4$};\node[above right] at (2,10) {$10$};
\node[below right] at (1.1,6.1) {$(1,7)$};
\draw[gray] (0,0+.08)--(0,0-.08);\node[below] at (0,0-.15) {$0$};
\draw[gray] (1,0+.08)--(1,0-.08);\node[below] at (1,0-.15) {$1$};
\draw[gray] (2,0+.08)--(2,0-.08);\node[below] at (2,0-.15) {$2$};
\end{scope}
\node at (6.8,.45) {$0,2$ 关于 $1$ 配对；$f(0)+f(2)=4+10=14$。};
\end{tikzpicture}

  \end{center}

  \medskip
  \textbf{\textcolor{CExample}{例 3（辨错：$b\ne0$ 就没有对称轴吗？）}}\par\nopagebreak[4]

  \textbf{题目：}$f'(x)=2x+2$ 的图像关于 $(0,2)$ 中心对称。判断原函数 $f(x)=x^2+2x$ 是否有对称轴。

  \textbf{解题步骤：}
  \begin{description}[style=nextline,leftmargin=2.8em,labelsep=0.5em,itemsep=0.3em]
    \item[\textbf{第一步（代入推广式）：}] $a=0,b=2$，所以 $f(t)-f(-t)=4t$。当 $t\ne0$ 时，一般不为零，故 $x=0$ 不是对称轴。
    \item[\textbf{第二步（直接看原函数）：}] $f(x)=x^2+2x=(x+1)^2-1$，显然关于 $x=-1$ 对称。
    \item[\textbf{第三步（得出正确判断）：}] $b\ne0$ 只能说明\textbf{不是关于 $x=a$ 对称}，不能断言没有任何对称轴。
  \end{description}
  \textbf{关键结论：}“$b\ne0$ 则没有对称轴”说法过强，应改成“$b\ne0$ 则 $x=a$ 不是对称轴”。
  \begin{center}
  \begin{tikzpicture}[font=\small,>=stealth,line join=round]
\draw[draw=teal!45!black,fill=green!3,rounded corners=3pt] (0,0) rectangle (13.6,5.9);
\node at (6.8,5.3500000000000005) {不是关于 $x=0$ 对称，仍然可能有另一条对称轴};
\begin{scope}[shift={(4.6,2.1)},x=1.1cm,y=0.75cm]
\draw[gray,->] (-3.2,0)--(1.7,0) node[right] {$x$};
\draw[gray,->] (0,-1.5)--(0,3.1);
\draw[teal!75!black,very thick,domain=-2.8:0.8,samples=100] plot(\x,{(\x)^2+2*\x});
\draw[teal!65!black,dashed] (-1,-1.25)--(-1,2.8);
\draw[red!65!black,dotted] (0,-1.25)--(0,2.8);
\fill[teal!75!black] (-1,-1) circle(2.5pt);
\fill[teal!75!black] (-2,0) circle(2.5pt);
\fill[teal!75!black] (0,0) circle(2.5pt);
\draw[gray] (-2,0+.08)--(-2,0-.08);\node[below] at (-2,0-.15) {$-2$};
\draw[gray] (-1,0+.08)--(-1,0-.08);\node[below] at (-1,0-.15) {$-1$};
\draw[gray] (0,0+.08)--(0,0-.08);\node[below] at (0,0-.15) {$0$};
\end{scope}
\node[anchor=west,align=left,text width=4.7cm] at (8,3.1) {$f(x)=(x+1)^2-1$。\\[10pt]绿虚线 $x=-1$：真正的轴。\\[10pt]红点线 $x=0$：不是轴。\\[10pt]$b\ne0$ 只排除指定的 $x=a$。};
\node at (6.8,.4) {$f^{\prime}(x)=2x+2$ 关于 $(0,2)$ 中心对称，但 $f$ 的轴在 $x=-1$。};
\end{tikzpicture}

  \end{center}

  \medskip
  \textbf{\textcolor{CExample}{例 4（高考型抽象条件：两个对称关系连用）}}\par\nopagebreak[4]

  \textbf{题目：}已知 $f$、$f'$ 的定义域均为 $\mathbb R$，$g(x)=f'(x)$，且 $f\bigl(\tfrac32-2x\bigr)$ 和 $g(2+x)$ 都是偶函数。证明 $f(x+2)=f(x)$。

  \textbf{解题步骤：}
  \begin{description}[style=nextline,leftmargin=2.8em,labelsep=0.5em,itemsep=0.3em]
    \item[\textbf{第一步（看第一组条件）：}] 由 $f(\tfrac32-2x)$ 为偶函数，得 $f(\tfrac32-2x)=f(\tfrac32+2x)$。因为 $x$ 可任意取值，故 $f$ 关于 $x=\tfrac32$ 对称，即 $f(3-u)=f(u)$。
    \item[\textbf{第二步（看第二组条件）：}] $g(2+x)$ 为偶函数，故 $g(2+x)=g(2-x)$，即 $f'$ 关于 $x=2$ 对称。由第二组规律，$f$ 关于 $(2,f(2))$ 中心对称：$f(4-u)=2f(2)-f(u)$。
    \item[\textbf{第三步（两次翻折）：}] 在第二个关系中令 $u=3-x$，得 $f(1+x)=2f(2)-f(3-x)$。由第一步 $f(3-x)=f(x)$，得到 $f(x+1)=2f(2)-f(x)$。
    \item[\textbf{第四步（再向右移一次）：}] $f(x+2)=2f(2)-f(x+1)=2f(2)-[2f(2)-f(x)]=\boxed{f(x)}$。
  \end{description}
  \textbf{关键结论：}两组对称合用，可以不用解析式就推出周期。图中选择满足条件的 $f(x)=\sin(\pi x)$ 展示变化；它不代替上面的一般证明。这里得到 $2$ 是\textbf{一个周期}，不一定是\textbf{最小正周期}。
  \begin{center}
  \begin{tikzpicture}[font=\small,>=stealth,line join=round]
\draw[draw=teal!45!black,fill=green!3,rounded corners=3pt] (0,0) rectangle (13.6,6.4);
\node at (6.8,5.8500000000000005) {取满足例题条件的 $f(x)=\sin(\pi x)$：一格反高，两格还原};
\begin{scope}[shift={(1.3,2.7)},x=1.75cm,y=1.15cm]
\draw[gray,->] (-0.1,0)--(4.15,0) node[right] {$x$};
\draw[gray,->] (0,-1.4)--(0,1.5);
\draw[teal!75!black,very thick,domain=0:4,samples=100] plot(\x,{sin(180*\x)});
\draw[teal!65!black,dashed] (1.5,-1.25)--(1.5,1.25);
\fill[blue!70!black] (2,0) circle(2.5pt);
\node[above] at (1.5,1.5) {轴 $x=\frac32$};
\node[below] at (2,-1.5) {中心 $(2,0)$};
\draw[gray] (0,0+.08)--(0,0-.08);\node[below] at (0,0-.15) {$0$};
\draw[gray] (1,0+.08)--(1,0-.08);\node[below] at (1,0-.15) {$1$};
\draw[gray] (2,0+.08)--(2,0-.08);\node[below] at (2,0-.15) {$2$};
\draw[gray] (3,0+.08)--(3,0-.08);\node[below] at (3,0-.15) {$3$};
\draw[gray] (4,0+.08)--(4,0-.08);\node[below] at (4,0-.15) {$4$};
\fill[orange!85!black] (0.25,0.7071067811865476) circle(2.5pt);
\fill[orange!85!black] (1.25,-0.7071067811865475) circle(2.5pt);
\fill[orange!85!black] (2.25,0.7071067811865474) circle(2.5pt);
\draw[orange!85!black,dotted] (.25,{sqrt(.5)})--(2.25,{sqrt(.5)});
\end{scope}
\node[anchor=west,align=left,text width=3.7cm] at (9.2,3.3) {$f(x+1)=-f(x)$\\（本例 $f(2)=0$）。\\[12pt]再走一格：\\$f(x+2)=f(x)$。\\[12pt]图像只是代表例子，\\一般证明见正文。};
\end{tikzpicture}

  \end{center}
\end{examplebox}
